\subsection{\texorpdfstring{\(\mathbf{E}'\) 的等度连续子集}{\textbackslash mathbf\{E\}\textquotesingle{} 的等度连续子集}}

在本节中，E 表示一个局部凸空间，\(E'\) 表示其对偶。每当我们谈到一个集 M 的极 \(M^{\circ}\)（M 位于 E 中，相应地，位于 \(E'\) 中）时，除非另有说明，我们总是指相对于 E 与 \(E'\) 之间的对偶而言的 M 的极。回忆：若 V 是 E 中 0 的一个闭凸均衡邻域，我们有 \(V^{\circ\circ}=V\)（II, 第 45 页, 推论 3）。

命题 7. --- 设 M 是 E' 的一个子集。下列条件等价：

\begin{enumerate}
\def\labelenumi{(\roman{enumi})}
\item
  M 是等度连续的；
\item
  \(\mathbf{M}\) 含于 E 中 0 的一个邻域的极中；
\item
  \(\mathbf{M}\) 的极是 E 中 0 的一个邻域。
\end{enumerate}

若 \(\mathbf{M}\) 等度连续，则存在 0 的一个凸均衡邻域 \(V\)，使得 \(|u(x)| \leqslant 1\) 对所有 \(x \in V\) 及所有 \(u \in M\) 成立；于是我们有 \(\mathbf{M} \subset V^{\circ}\)，从而 (i) 蕴含 (ii)。用同样的记号，若 \(\mathbf{M} \subset V^{\circ}\)，则 \(\mathbf{V} \subset V^{\circ \circ} \subset M^{\circ}\)，从而 (ii) 蕴含 (iii)。最后，若 \(\mathbf{M}^{\circ}\) 包含 0 的一个凸均衡邻域 \(V\)，则 \(\mathbf{M} \subset M^{\circ \circ} \subset V^{\circ}\)，且关系 \(x \in \varepsilon V\)、\(u \in M\) 蕴含 \(|u(x)| \leqslant \varepsilon\)（对所有 \(\varepsilon > 0\)），这就证明了 (iii) 蕴含 (i)。

我们注意到，每个 \(x \in E\) 定义了一个映射 \(j(x): u \mapsto u(x)\)，它从 \(E'\) 到 K 中。因此我们可以谈论 E 上的 S-拓扑，其中 S 是 \(E'\) 的子集的一个族：这是 \(K^{E'}\) 上的 S-拓扑在 j 下的原像。我们立即验证：若 S 是 \(E'\) 上的一个凸有界结构，则 S 中诸集的极构成 E 上 S-拓扑的 0 的一个基本邻域系。特别地，当 S 是 \(E'\) 的等度连续子集的族时即是如此，且命题 7 蕴含：

推论 1. --- E 的拓扑与在 E' 的等度连续子集上一致收敛的拓扑相同。

更一般地，设 F 是一个局部凸空间；每个 \(u \in \mathcal{L}(\mathrm{E}; \mathrm{F})\) 定义了一个映射 \(j(u):(x, f) \mapsto f(u(x))\)，它从 \(E \times F'\) 到 K 中（即到 R 或 C 中）。这使我们能够在空间 \(\mathcal{L}(\mathrm{E}; \mathrm{F})\) 上定义在 \(E \times F'\) 的子集所成之集上一致收敛的拓扑。特别地：

推论 2. --- 设 S 是 E 的有界子集的一个族。\(\mathcal{L}(\mathrm{E};\mathrm{F})\) 上的 S-拓扑是在形如 \(\mathbf{A}\times\mathbf{B}\subset\mathbf{E}\times\mathbf{F}'\) 的集合上一致收敛的拓扑，其中 A 属于 S，而 B 属于 F' 的等度连续子集的族。

对每个 \(u \in \mathcal{L}(\mathrm{E}; \mathrm{F})\)、每个 \(A \in S\) 以及 F 中 0 的每个闭凸均衡邻域 V，关系 \(u(\mathrm{A}) \subset \mathrm{V}\) 等价于 \(\ll j(u) (\mathrm{A} \times \mathrm{V}^{\circ})\) 含于 K 的单位球中 \(\gg\)。

命题 8. --- 设 H 是从 E 到一个局部凸空间 F 中的线性映射的一个族。为使 H 等度连续，必要且充分的是：对 F 的对偶 F' 中每个等度连续子集 X，线性型 \(f \circ u\)（\(f \in X\) 且 \(u \in H\)）所成之集是等度连续的。

显然该条件是必要的。假设它成立，并设 V 是 F 中 0 的一个闭凸均衡邻域。由于 \(V^{\circ}\) 等度连续，存在 E 中 0 的一个邻域 W，使得 \(\left|f(u(x))\right| \leqslant 1\) 对所有 \(x \in W\)、\(u \in H\) 及 \(f \in V^{\circ}\) 成立；换言之，\(u(W) \subset V^{\circ\circ} = V\) 对所有 \(u \in H\) 成立，故 H 等度连续。

